\section{Background and Related Works}\label{sec:related}

\subsection{Automatic Heuristic Design}
Automatic heuristic algorithm design is commonly known as hyper-heuristics~\cite{burke2013hyper,burke2019classification,stutzle2019automated}. With various effective methodologies~\cite{blot2016mo,lopez2016irace,akiba2019optuna} and frameworks~\cite{burke2019classification}, one can tune heuristics or combine different algorithmic components in an automatic manner. Much effort has been made to use machine learning techniques in automatic algorithm design~\cite{bengio2021machine,chen2022learning,he2021automl,li2023survey}. Among them, genetic programming~\cite{mei2022explainable,jia2022learning} provides an explainable approach to algorithm design. However, it requires hand-crafted algorithmic components and domain knowledge.

\subsection{LLMs for Heuristic Design}
Over the last few years, the ability of large language models has increased significantly~\cite{naveed2023comprehensive}. Recently, some effort has been made to use LLMs as basic algorithmic components to improve the performance of algorithms~\cite{yang2023large,guo2023towards}. Most of these works adopt LLM as optimizers~\cite{yang2023large} to directly generate new trial solutions through in-context learning. This approach faces challenges when applied to complex problems with large search space~\cite{yang2023large,nasir2023llmatic,zhao2023large,liu2023large}. Others integrate LLMs to assist algorithm design to extract deep algorithm features for heuristic selection~\cite{wu2023llm}, provide a guide for heuristic~\cite{shah2023navigation}, and design an algorithmic component~\cite{xiao2023llm}.  However, designing a competitive heuristic is still a challenge for standalone LLMs with prompt engineering.

\subsection{LLMs + EC}
Evolutionary computation is a generic optimization principle inspired by natural evolution~\cite{back1997handbook,eiben2015evolutionary}. Integration of EC in the prompt engineering of LLMs is very promising in improving performance in various domains~\cite{guo2023connecting,lehman2023evolution,wu2024evolutionary}. Evolutionary methods have been adopted in both code generation~\cite{liventsev2023fully,ma2023eureka,Lehman2024,hemberg2024evolving} and text generation~\cite{guo2023connecting,fernando2023promptbreeder,xu2023wizardlm}. The most related work to our effort is FunSearch~\cite{romera2024mathematical}, an evolutionary framework with LLMs to search functions automatically. Algorithms generated by FunSearch outperform hard-crafted algorithms on some optimization problems. However, FunSearch is computationally expensive and usually needs to generate millions of programs (i.e., queries to LLMs) to identify an effective heuristic function, which is not very practical for many users.


\vspace{1pt}